% 
%  chapter2.tex
%  
%  Created by Matilde Pós-de-Mina Pato on 2012/10/09.
%
\chapter{Conclusões}
\label{cha:Conclusion}

\section{...}
Qualquer trabalho deverá ser escrito como um artigo, i.e. a linguagem deve ser clara, objectiva, escrita em discurso directo e com frases curtas \cite{gustavii2016write, glasman2010science}.

A redação deve ser feita com frases curtas e objectivas, organizadas de acordo com a estrutura do trabalho, dando destaque a cada uma das partes abordadas, assim apresentadas: Introdução - Informar, em poucas palavras, o contexto em que o trabalho se insere, sintetizando a problemática estudada. Objetivo - Deve ser explicitado claramente. Métodos - Destacar os procedimentos metodológicos adoptados. Resultados - Destacar os mais relevantes para os objectivos pretendidos. Os trabalhos de natureza quantitativa devem apresentar resultados numéricos, assim como seu significado estatístico. Conclusões - Destacar as conclusões mais relevantes, os estudos adicionais recomendados e os pontos positivos e negativos que poderão influir no conhecimento. 

% ---------

\section{...}
Qualquer trabalho deverá ser escrito como um artigo, i.e. a linguagem deve ser clara, objectiva, escrita em discurso directo e com frases curtas \cite{gustavii2016write, glasman2010science}.

A redação deve ser feita com frases curtas e objectivas, organizadas de acordo com a estrutura do trabalho, dando destaque a cada uma das partes abordadas, assim apresentadas: Introdução - Informar, em poucas palavras, o contexto em que o trabalho se insere, sintetizando a problemática estudada. Objetivo - Deve ser explicitado claramente. Métodos - Destacar os procedimentos metodológicos adoptados. Resultados - Destacar os mais relevantes para os objectivos pretendidos. Os trabalhos de natureza quantitativa devem apresentar resultados numéricos, assim como seu significado estatístico. Conclusões - Destacar as conclusões mais relevantes, os estudos adicionais recomendados e os pontos positivos e negativos que poderão influir no conhecimento.

